%% =====================================================================
%%  CHAPTER 6: DESIGN ALTERNATIVES, SELECTION, INNOVATION   (budget 8 pages)
%%  chapters/ch6-design.tex
%%  Every Pugh score c_ij is derived from an engineering metric (rules in
%%  Table 6.3); model parameters are labelled as planning assumptions.
%% =====================================================================
\pcChapter{alt}{8}{חלופות תכן, בחירת פתרון וחדשנות}
\begin{pcguide}[מה כותבים כאן]
\latin{2}--\latin{3} חלופות תכן והשוואה מנומקת, כדי שהבחירה תהיה מוצדקת. אפשר לשלב שני פתרונות קיימים;
אפשר לכלול פתרון יקר או לא זמין. סמנו איזו חלופה היא הפתרון המוצע בפרק~\ref{ch:solution}.
תקציב: \latin{2.5} + \latin{2} + \latin{2.5} + \latin{1} עמודים לפי הסעיפים.
\end{pcguide}

\section{חלופות תכן}\label{sec:alts}
נבחנו שלוש חלופות שונות עקרונית (טבלה~\ref{tab:alts}). \textbf{חלופה א --- בסיס ידני:} ראיונות מומחים, קובצי
\latin{Word/Excel/PowerPoint}, חפיפה אישית ובחינה ידנית. \textbf{חלופה ב --- מאגר \latin{LMS} וחיפוש מילולי:} תיעוד דיגיטלי
מרכזי וחיפוש מילולי \latin{(BM25)}, סרטונים ומבחנים שכותב מדריך, ללא הפקה מבוססת \latin{RAG}. \textbf{חלופה ג --- \latin{OrbiTech}:}
קליטת מסמכים ולקחים, אחזור היברידי מבוסס הרשאות, מודל שפה מקומי להפקת טיוטות, שרשרת מקורות, אישור מומחה, מבחנים ותבניות
\threeD{} פרמטריות שמאוכלסות מתצורות \latin{JSON} שנוצרות אוטומטית; מסלול \twoD{} חלופי משמש לנגישות ולהשוואה.
\begin{table}[!htbp]
\centering
\caption{השוואת חלופות התכן}\label{tab:alts}
\begin{pctable}{H{0.22},H{0.26},H{0.26},H{0.26}}
  \pchead{היבט & חלופה א: ידני & חלופה ב: \latin{LMS} וחיפוש & חלופה ג: \latin{OrbiTech}}
  \pcrow{עיקרון פעולה & ראיונות, קבצים משרדיים וחפיפה אישית & מאגר דיגיטלי, חיפוש מילולי, סרטונים ומבחנים ידניים & \latin{Grounded RAG} מקומי, אישור מומחה ותרגול \latin{3D} מתצורת \latin{JSON}}
  \pcrow{מקום פריסה ותלות חיצונית & מקומי; ללא תלות & מקומי או שירות ענן, לפי המערכת & מקומי ומבודד; ללא רשת יוצאת}
  \pcrow{יתרונות & סיכון טכנולוגי נמוך; שליטה אנושית באישור & פשטות וניהול למידה; מסייע לניהול גרסאות & הפקה מהירה, עקיבות למקור, עדכון אוטומטי ומדידת כשירות}
  \pcrow{חסרונות ומגבלות & זמן הכנה ארוך; מעקב גרסאות ידני & עומס יצירת התוכן נשאר על המדריך & מורכבות הטמעה, צורך ב-\latin{GPU}; עמידה באבטחה טרם הוכחה}
  \pcrow{עלות ומורכבות & שעות מומחה ומדריך & רישוי ותחזוקת מאגר & תשתית \latin{GPU}, זמן מומחה לאישור ותחזוקה}
  \pcrow{קשר לפתרון המוצע (פרק~\ref{ch:solution}) & בסיס ההשוואה ב-\hyp{2} & בסיס עקרוני לניהול למידה & \textbf{החלופה הנבחרת}}
\end{pctable}
\end{table}

\section{מטריצת החלטה משוקללת (\latin{Pugh})}\label{sec:pugh}
סכום הניקוד המשוקלל של חלופה \mth{j} מחושב לפי
\begin{equation}
  S_j=\sum_{i=1}^{n} w_i\, c_{ij},\qquad \sum_{i=1}^{n} w_i = 1,
\end{equation}
כאשר \mth{w_i} הוא משקל הקריטריון \mth{i} ו-\mth{c_{ij}} הציון של חלופה \mth{j} בסולם \rng{1}{5} (גבוה עדיף). כדי שהציונים לא
יהיו הערכות שרירותיות, כל \mth{c_{ij}} נגזר ממדד הנדסי מוגדר בכלל המרה קבוע מראש (טבלה~\ref{tab:rules}), ומהערך המחושב של המדד
במודל הנדסי שפרמטריו מפורטים בהמשך. פרמטרי המודל הם \textbf{הנחות תכנון} שיימדדו בפיילוט (\latin{D4} ושבוע \latin{5}), ולא
תוצאות; טבלה~\ref{tab:pugh} מציגה לכל תא את ערך המדד ואת הציון שנגזר ממנו.
\begin{table}[!htbp]
\centering
\caption{מטריצת החלטה משוקללת: ערך המדד ההנדסי \mth{\rightarrow} הציון \mth{c_{ij}}; סכום המשקלים \latin{1}}\label{tab:pugh}
\pughReset
\begin{pctable}{H{0.33},C{0.09},C{0.19},C{0.19},C{0.20}}
  \pchead{קריטריון (מדד) & משקל \mth{w_i} & חלופה א & חלופה ב & חלופה ג}
  \pughRowEv{חשיפה להזיה: שיעור טענות לא נתמכות (\latin{\%})}{0.20}{5}{0}{5}{0}{4}{2.4}
  \pughRowEv{אבטחה ושליטה מקומית: בקשות לרשת חוץ}{0.20}{5}{0}{3}{SaaS}{4}{0 (design)}
  \pughRowEv{מאמץ הפקה יחסי \mth{E} (ידני \mth{=1}) }{0.15}{2}{1.00}{3}{0.82}{4}{0.52}
  \pughRowEv{נטל תשתית: \latin{VRAM} (\latin{GB})}{0.15}{5}{0}{4}{0}{3}{9.4}
  \pughRowEv{זמן הפצת עדכון גרסה \mth{T_{upd}} (שעות)}{0.10}{2}{24}{3}{16}{5}{1}
  \pughRowEv{ערוצי העברת ידע נתמכים (\latin{0}--\latin{3})}{0.20}{3}{2}{2}{1}{4}{3}
  \pughTotalsAbs
\end{pctable}
\pughCheck
\end{table}
\begin{table}[!htbp]
\centering
\caption{כללי ההמרה ממדד הנדסי לציון, ונימוק המשקולות}\label{tab:rules}
\begin{pctable}{H{0.22},C{0.08},H{0.28},H{0.20},H{0.22}}
  \pchead{קריטריון & \mth{w_i} & סף לציונים \latin{5/4/3/2/1} & מקור המדד & נימוק המשקל}
  \pcrow{חשיפה להזיה & \num{0.20} & \latin{$\le$1 / $\le$5 / $\le$10 / $\le$20 / $>$20 (\%)} & מודל שכבות בקרה (להלן) \cite{ji2023survey,barnett2024seven} & פלט לא מעוגן הוא הסיכון החמור ביותר (\risk{03}); \req{03}}
  \pcrow{אבטחה ושליטה מקומית & \num{0.20} & \latin{0} מאומת / \latin{0} בתכן / אחסון עצמי / \latin{SaaS} בלבד / ענן & \req{09}, \req{10} \cite{zeng2024good} & סביבה ביטחונית-תעשייתית}
  \pcrow{מאמץ הפקה \mth{E} & \num{0.15} & \latin{$\le$0.5 / $\le$0.7 / $\le$0.9 / $\le$1.0 / $>$1.0} & מודל מאמץ (להלן); יעד \req{04} & רכיב \mth{T_{prep}}, \hyp{2}}
  \pcrow{נטל תשתית (\latin{VRAM}) & \num{0.15} & ללא \latin{GPU} / ללא \latin{GPU}, שילוב / \latin{$\le$16 / $\le$24 / $>$24 (GB)} & חישוב זיכרון (להלן) & צוות של שניים, \latin{14} שבועות (טבלה~\ref{tab:constraints})}
  \pcrow{הפצת עדכון \mth{T_{upd}} & \num{0.10} & \latin{$\le$2 / $\le$8 / $\le$16 / $\le$40 / $>$40 (h)} & יעד \req{11}: עד יום עבודה & ערבוב גרסאות \cite{zheng2025knowshift}}
  \pcrow{ערוצי העברת ידע & \num{0.20} & \latin{3} מאומת / \latin{3 / 2 / 1 / 0} & מסמך, ראיון סמוי, תרגול החלטה \cite{bodding2025mr} & רכיבי שאלת המחקר: \req{16}, \hyp{4}}
\end{pctable}
\end{table}
\textbf{מודל החשיפה להזיה.} הזיה נשארת בתוצר המאושר רק אם עברה שלוש שכבות בקרה: מודל השפה מייצר שיעור גולמי \mth{U_{gen}} של טענות
לא נתמכות; מנגנון ההימנעות והאימות מסנן חלק \mth{d_{abs}} מהן; ומומחה מזהה חלק \mth{d_{SME}} מהשאר:
\begin{equation}
  U_{final}=U_{gen}\,(1-d_{abs})\,(1-d_{SME})=0.20\times0.60\times0.20=0.024\ (2.4\%).
\end{equation}
הערכים \mth{U_{gen}=20\%}, \mth{d_{abs}=40\%} ו-\mth{d_{SME}=80\%} הם הנחות תכנון; אם יעילות המומחה תרד ל-\pct{50}, \mth{U_{final}}
יעלה ל-\pct{6} והציון ירד ל-\latin{3}. בחלופות א ו-ב אין שלב יצירה, ולכן אין הזיה, אך נשאר סיכון לטעות אנוש שאינו נמדד. \textbf{מודל
המאמץ.} המאמץ ליחידה מאושרת הוא \mth{E=\sum_k s_k(1-\Delta_k)}, כאשר \mth{s_k} חלק השלב בבסיס הידני (איסוף \mth{0.30},
כתיבה \mth{0.50}, בקרה \mth{0.20}; הנחה שתיבדק ב-\latin{D4}) ו-\mth{\Delta_k} השינוי היחסי בשלב. בחלופה ב החיפוש מקצר את האיסוף ב-\pct{60}:
\mth{E_B=0.12+0.50+0.20=0.82}. בחלופה ג הטיוטה מקצרת את הכתיבה ב-\pct{80}, והבקרה מתארכת ב-\pct{50} משום שהמומחה בודק
ציטוטים: \mth{E_C=0.12+0.10+0.30=0.52}, כלומר קרוב ליעד \req{04} (\mth{\le0.50}) אך מעליו --- היעד הדוק, ויבדק בפיילוט.
\begin{figure}[!htbp]
\centering
\pcFig{pugh_bar}{%
\pcAutoFit{\begin{pctikz}
\begin{axis}[ybar stacked,width=10.5cm,height=6.2cm,symbolic x coords={A,B,C},xtick=data,
  xticklabels={\tn{חלופה א},\tn{חלופה ב},\tn{חלופה ג}},ymin=0,ymax=4.6,bar width=24pt,
  ymajorgrids,ylabel={\tn{ניקוד משוקלל}},enlarge x limits=0.25,
  legend style={at={(0.5,-0.18)},anchor=north,legend columns=2,font=\scriptsize,draw=none},
  cycle list={{fill=pcPrimary,draw=white},{fill=pcGood,draw=white},{fill=pcWarn,draw=white},{fill=pcBad,draw=white},{fill=pcGrey,draw=white},{fill=pcDark,draw=white}}]
\pughSeries
\end{axis}
\end{pctikz}}}
\caption{פירוק הניקוד המשוקלל של כל חלופה לתרומות הקריטריונים (מחושב מטבלה~\ref{tab:pugh}; ערכי המדדים הם פרמטרי תכנון)}\label{fig:pughbar}
\end{figure}
\textbf{ניתוח רגישות.} משנים משקל של קריטריון אחד, \mth{w_k \rightarrow w_k'}, ומנרמלים את האחרים כך שהסכום יישאר \latin{1}:
\begin{equation}
  w_i' = w_i\,\frac{1-w_k'}{1-w_k}\quad (i\neq k),\qquad S_j' = \sum_{i} w_i'\, c_{ij}.
\end{equation}
טבלה~\ref{tab:sens} מראה את התוצאה. הבחירה בחלופה ג יציבה כשהמשקל הגדול הוא על מאמץ ההפקה, על העברת הידע או על משקולות שוות, אך
לא כשהמשקל הגדול הוא על נאמנות, אבטחה או קלות הטמעה: אז החלופה הידנית עולה לראש. נקודות המעבר: חלופה א עוקפת את ג אם
\mth{w_1\ge0.275} (נאמנות), \mth{w_2\ge0.275} (אבטחה) או \mth{w_4\ge0.195} (נטל תשתית, לעומת \mth{0.15} בבסיס). זו תוצאה הגיונית
ולא פגם: חלופה ג נבחרת משום שהיא משרתת את מטרות הפרויקט (מאמץ הפקה, עדכון והעברת ידע), ובמחיר סיכון ותשתית שהפיילוט נועד לבדוק.
\begin{table}[!htbp]
\centering
\caption{ניתוח רגישות: האם הבחירה יציבה לשינוי משקולות?}\label{tab:sens}
\begin{pctable}{H{0.24},C{0.26},C{0.09},C{0.09},C{0.09},C{0.14}}
  \pchead{תרחיש (המשקל המשתנה) & \mth{(w_1,\dots,w_6)} & \mth{S_A} & \mth{S_B} & \mth{S_C} & החלופה המובילה}
  \pcrow{בסיס & \numlist{(.20,.20,.15,.15,.10,.20)} & \num{3.85} & \num{3.35} & \num{3.95} & ג}
  \pcrow{נאמנות דומיננטית (\mth{w_1=0.35}) & \numlist{(.35,.16,.12,.12,.08,.16)} & \num{4.07} & \num{3.66} & \num{3.96} & א}
  \pcrow{אבטחה דומיננטית (\mth{w_2=0.35}) & \numlist{(.16,.35,.12,.12,.08,.16)} & \num{4.07} & \num{3.28} & \num{3.96} & א}
  \pcrow{נטל תשתית דומיננטי (\mth{w_4=0.35}) & \numlist{(.15,.15,.12,.35,.08,.15)} & \num{4.12} & \num{3.50} & \num{3.73} & א}
  \pcrow{מאמץ הפקה דומיננטי (\mth{w_3=0.30}) & \numlist{(.17,.17,.30,.12,.08,.17)} & \num{3.52} & \num{3.29} & \num{3.96} & ג}
  \pcrow{משקולות שוות & \numlist{(.17,.17,.17,.17,.17,.17)} & \num{3.67} & \num{3.33} & \num{4.00} & ג}
\end{pctable}
\end{table}
\textbf{נימוק הבחירה והשער.} הפרש של \num{0.10} בתרחיש הבסיס אינו מכריע, ולכן הבחירה בחלופה ג מותנית בשערים שנקבעו מראש: נאמנות
מדודה של לפחות \pct{90}, אפס בקשות לרשת חוץ, וקצב יצירה העומד ביעד \req{08} (להלן). אם אחד מהם ייכשל, הפרויקט יעבור
לגיבוי: חלופה א משופרת באחזור \latin{RAG} בלבד, כפי שמגדיר מאמר הדגל (פרק~\ref{ch:lit}) וללא שכבת ה-\latin{3D} המתקדמת.

\section{ארכיטקטורת המערכת ותרשים זרימת נתונים}\label{sec:arch}
תרשים זרימת הנתונים (איור~\ref{fig:dfd}) כולל שלוש ישויות אנושיות: מנהל ידע, מומחה תוכן מורשה ולומד. שני מאגרי נתונים מרכזיים:
מאגר מקטעים בגרסאות, ויומן מקורות, אישורים ופעולות. מנוע ההפקה אינו רשאי לגשת למסמך שאינו מותר לתפקיד, וניהול הגרסאות מתבצע גם לפני
ההפקה וגם במועד שליפת הפריט. פרטי הטכנולוגיה (דגם, אחזור, מסד וקטורי) מפורטים בפרק~\ref{ch:solution}; כאן מפורטים הסכמות, צינור הקליטה והגרסאות,
תהליך האישור האנושי ומכונת המצבים של הסימולציה.
\begin{figure}[!htbp]
\centering
\pcFig{dfd}{%
\pcAutoFit{\begin{pctikz}[x=1cm,y=1cm]
  \node[pcentity,text width=2.2cm] (e1) at (0,0) {\tn{מנהל ידע: מסמכים ולקחים מאושרים}};
  \node[pcproc] (p1) at (-3.2,0) {\tn{\latin{1.} פענוח ואינדוקס}};
  \node[pcproc] (p2) at (-6.4,0) {\tn{\latin{2.} אחזור \latin{RAG} והפקה עם סימוכין}};
  \node[pcproc] (p3) at (-9.6,0) {\tn{\latin{3.} אישור מומחה \latin{(SME)}}};
  \node[pcentity,text width=2.2cm] (e2) at (-12.8,0) {\tn{לומד: יחידה, מבחן ותרחיש \latin{3D}}};
  \node[pcstore] (d1) at (-4.8,-3.1) {\tn{מאגר מקטעים: גרסה, עמוד, הרשאה}};
  \node[pcstore] (d2) at (-11.2,-3.1) {\tn{יומן מקורות, אישורים ופעולות}};
  \node[pcentity,text width=2.2cm] (e3) at (-9.6,3.0) {\tn{מומחה תוכן מורשה}};
  \draw[pcflow] (e1) -- node[above,font=\scriptsize]{\tn{מסמכים}} (p1);
  \draw[pcflow] (p1) -- node[above,font=\scriptsize]{\tn{שאילתה}} (p2);
  \draw[pcflow] (p2) -- node[above,font=\scriptsize]{\tn{טיוטה}} (p3);
  \draw[pcflow] (p3) -- node[above,font=\scriptsize]{\tn{פריט מאושר}} (e2);
  \draw[pcflow] (p1) -- node[right,font=\scriptsize]{\tn{כתיבה}} (d1);
  \draw[pcflow] (d1) -| node[near start,above,font=\scriptsize]{\tn{ראיות}} (p2);
  \draw[pcflow] (p3) -- node[left,font=\scriptsize]{\tn{רישום}} (d2);
  \draw[pcflow] (e2) |- node[near start,left,font=\scriptsize]{\tn{פעולות \latin{3D}}} (d2);
  \draw[pcflow] ($(e3.south)+(0.35,0)$) -- node[right,font=\scriptsize]{\tn{אישור}} ($(p3.north)+(0.35,0)$);
  \draw[pcflow] ($(p3.north)+(-0.35,0)$) -- node[left,font=\scriptsize]{\tn{לבדיקה}} ($(e3.south)+(-0.35,0)$);
\end{pctikz}}}
\caption{תרשים זרימת נתונים \latin{(DFD)}: מקורות מאושרים, \latin{RAG}, בקרת מומחה, תרחישי \latin{3D} ומעקב למידה}\label{fig:dfd}
\end{figure}

\subsection{סכמות הנתונים}
הנתונים מאורגנים בחמש ישויות (טבלה~\ref{tab:schema}). המפתח \texttt{\latin{revision}} מופיע בכל ישות הנגזרת ממסמך, ולכן כל טענה, יחידה
וצומת תרחיש ניתנים לשחזור עד גרסת המקור שלהם; \texttt{\latin{status}} קובע אם הפריט יכול להגיע ללומד.
\begin{table}[!htbp]
\centering
\caption{סכמות הנתונים המרכזיות (שמות שדות באנגלית כפי שיישמרו במסד)}\label{tab:schema}
\begin{pctable}{H{0.19},H{0.46},H{0.35}}
  \pchead{ישות & שדות עיקריים & תפקיד}
  \pcrow{\texttt{\latin{Document}} & \texttt{\latin{document\_id, title, owner, revision, effective\_date, sha256, access\_level, status}} & מקור וגרסה; \texttt{\latin{sha256}} מזהה שינוי בקובץ}
  \pcrow{\texttt{\latin{Chunk}} & \texttt{\latin{chunk\_id, document\_id, revision, page, section, text, embedding\_id, valid\_from, valid\_to}} & יחידת האחזור; \texttt{\latin{valid\_to}} נסגר בעת החלפת גרסה}
  \pcrow{\texttt{\latin{Claim}} & \texttt{\latin{claim\_id, text, chunk\_ids[], origin (document|interview), status (draft|verified|stale|rejected)}} & טענה בודדת עם עיגון; מקור הריאיון נבדק מול מקור תקף}
  \pcrow{\texttt{\latin{Unit}} & \texttt{\latin{unit\_id, claim\_ids[], objectives[], questions[], status (draft, in\_review, approved, published, suspended), approved\_by, approved\_at}} & יחידת הדרכה ומבחן; פרסום רק במצב \texttt{\latin{approved}}}
  \pcrow{\texttt{\latin{ScenarioConfig}} & \texttt{\latin{scenario\_id, unit\_id, template\_id, scene\_version, nodes[\{node\_id, prompt, options[], correct, evidence\_id\}]}} & תצורת \latin{JSON} שמאכלסת תבנית \latin{3D} (פרק~\ref{ch:solution})}
  \pcrow{\texttt{\latin{EventLog}} & \texttt{\latin{scenario\_id, scene\_version, node\_id, evidence\_id, choice, is\_correct, response\_ms, timestamp, pseudo\_id}} & מדדי \hyp{4} ו-\req{14}; מזהה אקראי בלבד}
\end{pctable}
\end{table}

\subsection{צינור הקליטה וניהול הגרסאות}
איור~\ref{fig:pipeline} מציג את צינור הקליטה. כל קובץ שנקלט מקבל גיבוב \texttt{\latin{sha256}}; אם הגיבוב שונה מזה של הגרסה האחרונה
נפתחת גרסה חדשה. המקטעים של הגרסה הישנה אינם נמחקים: \texttt{\latin{valid\_to}} נקבע להיום, והם מסומנים \texttt{\latin{superseded}}. טענות
ההפניות אליהם הופכות ל-\texttt{\latin{stale}}, ויחידות המסתמכות עליהן עוברות למצב \texttt{\latin{suspended}} (\req{11}); בו בזמן נוצרת
בקשת ייצור מחודש לתור המומחה. כך עדכון נוהל אינו דורש חיפוש ידני של כל חומר ההדרכה שהושפע ממנו.
\begin{figure}[!htbp]
\centering
\pcFig{pipeline}{%
\pcAutoFit{\begin{pctikz}[x=1cm,y=1cm]
  \tikzset{st/.style={pcnode,text width=2.1cm,font=\scriptsize\sffamily}}
  \node[st] (s1) at (0,0) {\tn{1. בדיקת סיווג והרשאה; חישוב \latin{sha256}}};
  \node[st] (s2) at (-2.6,0) {\tn{2. פענוח, פירוק למקטעים (עמוד, סעיף)}};
  \node[st] (s3) at (-5.2,0) {\tn{3. השוואה לגרסה קודמת; פתיחת גרסה}};
  \node[st] (s4) at (-7.8,0) {\tn{4. אינדוקס: \latin{BM25} + וקטורי}};
  \node[st] (s5) at (-10.4,0) {\tn{5. סימון \latin{stale} והשהיית יחידות תלויות}};
  \node[st,fill=pcWarn!15,draw=pcWarn] (s6) at (-10.4,-2.0) {\tn{6. בקשת ייצור מחודש לתור המומחה}};
  \foreach \a/\b in {s1/s2,s2/s3,s3/s4,s4/s5}{\draw[pcflow] (\a) -- (\b);}
  \draw[pcflow] (s5) -- (s6);
\end{pctikz}}}
\caption{צינור הקליטה והגרסאות: שינוי גרסה משעה יחידות תלויות ופותח בקשת אישור מחודש}\label{fig:pipeline}
\end{figure}

\subsection{תהליך אישור המומחה \latin{(HITL)}}
מצבי הפריט ומעבריהם מוצגים באיור~\ref{fig:hitl}. טיוטה נבדקת אוטומטית (עיגון והימנעות), עוברת לתור המומחה, ורק פריט שאושר מתפרסם. פריט שנדחה
חוזר לטיוטה עם הערות, ושינוי גרסת מקור משעה פריט שפורסם ומחזיר אותו לבדיקה. קיבולת המומחה היא צוואר הבקבוק של התהליך; כדי לכמת אותה
מדמים את התור כמערכת \latin{M/M/1} (שיטה מקורס חקר ביצועים): בקצב הגעה \mth{\lambda} יחידות לשבוע וקצב בדיקה \mth{\mu},
\begin{equation}
  \rho=\frac{\lambda}{\mu},\qquad W_q=\frac{\rho}{\mu-\lambda}.
\end{equation}
בפיילוט מתוכננות \mth{\lambda=10} יחידות לשבוע; אם המומחה בודק \mth{\mu=14} יחידות, \mth{\rho=0.71} וזמן ההמתנה הממוצע
\mth{W_q=0.18} שבוע; ב-\mth{\mu=12} הוא עולה ל-\mth{0.42} שבוע (כיומיים), וב-\mth{\mu=11} ל-\mth{0.91} שבוע. כלומר ירידה של \pct{21} בזמינות המומחה (מ-\latin{14}
ל-\latin{11}) מגדילה את ההמתנה פי חמישה; זהו הבסיס לסיכון זמינות המומחים (\risk{06}) ולהגבלת קצב ההפקה ל-\latin{10} יחידות
בשבוע.
\begin{figure}[!htbp]
\centering
\pcFig{hitl}{%
\pcAutoFit{\begin{pctikz}[x=1cm,y=1cm]
  \tikzset{ms/.style={pcnode,text width=1.9cm,font=\scriptsize\sffamily,minimum height=7mm}}
  \node[ms] (m1) at (0,0) {\tn{טיוטה \latin{(draft)}}};
  \node[ms] (m2) at (-2.7,0) {\tn{נבדק אוטומטית}};
  \node[ms] (m3) at (-5.4,0) {\tn{בבדיקת מומחה \latin{(in\_review)}}};
  \node[ms,fill=pcGood!15,draw=pcGood] (m4) at (-8.1,0) {\tn{מאושר \latin{(approved)}}};
  \node[ms,fill=pcGood!25,draw=pcGood] (m5) at (-10.8,0) {\tn{מפורסם \latin{(published)}}};
  \node[ms,fill=pcBad!10,draw=pcBad] (m6) at (-5.4,-1.8) {\tn{נדחה \latin{(rejected)}}};
  \node[ms,fill=pcWarn!15,draw=pcWarn] (m7) at (-10.8,-1.8) {\tn{מושהה \latin{(suspended)}}};
  \foreach \a/\b in {m1/m2,m2/m3,m3/m4,m4/m5}{\draw[pcflow] (\a) -- (\b);}
  \draw[pcflow] (m3) -- (m6);
  \draw[pcflow] (m6.east) -| node[near start,above,font=\scriptsize]{\tn{תיקון}} (m1.south);
  \draw[pcflow] (m5) -- node[right,font=\scriptsize]{\tn{שינוי גרסה}} (m7);
  \draw[pcflow] (m7.east) -- node[above,font=\scriptsize]{\tn{לבדיקה מחדש}} (m6.west |- m7.east) -- (m3.south west);
\end{pctikz}}}
\caption{מכונת המצבים של אישור פריט: רק מצב \texttt{\latin{approved}} מגיע ללומד; שינוי גרסה מחזיר פריט לבדיקה}\label{fig:hitl}
\end{figure}

\subsection{מכונת המצבים של סימולציית \latin{3D}}
הסימולציה אינה סצנה שנבנתה ידנית אלא תבנית פרמטרית שטוענת תצורת \texttt{\latin{ScenarioConfig}}. בכל צומת התבנית מציגה את המצב, הלומד בוחר, והמערכת
משווה לתשובה הנכונה בתצורה, מציגה משוב מעוגן (\texttt{\latin{evidence\_id}}) ורושמת אירוע. איור~\ref{fig:sim} מציג את המצבים; תנאי ההצלחה של כל משימה
מוגדרים בתצורה, ומהם נגזר שיעור ההשלמה (\req{15}).
\begin{figure}[!htbp]
\centering
\pcFig{sim_fsm}{%
\pcAutoFit{\begin{pctikz}[x=1cm,y=1cm]
  \tikzset{ms/.style={pcnode,text width=2.0cm,font=\scriptsize\sffamily,minimum height=7mm}}
  \node[ms,fill=pcDark,text=white] (t0) at (0,0) {\tn{טעינת תצורת \latin{JSON}}};
  \node[ms] (t1) at (-2.9,0) {\tn{תצפית בסצנה}};
  \node[ms] (t2) at (-5.8,0) {\tn{בחירה בצומת החלטה}};
  \node[ms] (t3) at (-8.7,0) {\tn{משוב מעוגן + רישום אירוע}};
  \node[ms,fill=pcGood!20,draw=pcGood] (t4) at (-11.6,0) {\tn{סיום: חישוב השלמה}};
  \foreach \a/\b in {t0/t1,t1/t2,t2/t3,t3/t4}{\draw[pcflow] (\a) -- (\b);}
  \draw[pcflow] (t3.south) -- ++(0,-0.9) -| node[near start,below,font=\scriptsize]{\tn{צומת הבא}} (t1.south);
\end{pctikz}}}
\caption{מכונת המצבים של הסימולציה: כל החלטה נרשמת עם מזהה ראיה לפני המעבר לצומת הבא}\label{fig:sim}
\end{figure}
\textbf{תקציב \latin{VRAM} וקצב.} זיכרון הכרטיס לדגם \mth{P} פרמטרים בדיוק של \mth{b} סיביות הוא
\begin{equation}
  M_{VRAM}\approx \frac{P\cdot b}{8}+2\,L\,H_{kv}\,d_h\,B\,N_{ctx}+M_{other},
\end{equation}
כאשר \mth{L} מספר השכבות, \mth{H_{kv}} מספר ראשי המפתח-ערך, \mth{d_h} ממד הראש, \mth{B} מספר הבתים לאלמנט ו-\mth{N_{ctx}} אורך ההקשר. עבור דגם
בגודל \latin{8B} בדיוק \latin{4} סיביות עם \mth{L=32}, \mth{H_{kv}=8}, \mth{d_h=128} ו-\mth{B=2} (תצורה אופיינית לדגמי משפחה זו):
משקלים \mth{4.0} \latin{GB}, זיכרון הקשר \mth{1.1} \latin{GB} ל-\mth{8192} טוקנים (\mth{4.3} \latin{GB} ל-\latin{32,768}),
מודל ייצוג ומדרג כ-\mth{1.3} \latin{GB}, תקורת הפעלה כ-\mth{1.0} \latin{GB} ותצוגת \latin{3D} כ-\mth{2.0} \latin{GB}: סך הכול
כ-\mth{9.4} \latin{GB} (טבלה~\ref{tab:budget}), ולכן מומלץ כרטיס בעל זיכרון של יותר מ-\latin{16} \latin{GB} (למשל \latin{24} \latin{GB}), כדי להשאיר מרווח להקשר ארוך ולהרצה מקבילה של המודל והתצוגה.
זמן התגובה לשאילתה: \mth{T=T_{retr}+T_{pre}+N_{out}/v}; ב-\mth{T_{retr}=0.3} \latin{s}, \mth{T_{pre}=1.5} \latin{s} ו-\mth{N_{out}=300} טוקנים,
יעד \req{08} (\mth{T\le10} \latin{s}) מחייב קצב \mth{v\ge37} טוקנים לשנייה; בקצב \mth{40} נקבל \mth{9.3} \latin{s}, ובקצב \mth{20} --- \mth{16.8} \latin{s}.
הקצב ייבדק בשבוע \latin{5}, ואם יהיה נמוך מהסף ייעשה מעבר לדגם קטן יותר או קיצור הפלט.
\begin{table}[!htbp]
\centering
\caption{תקציב הנדסי מחושב: זיכרון \latin{GPU} וזמן תגובה (פרמטרי תכנון; יימדדו בשבוע \latin{5})}\label{tab:budget}
\begin{pctable}{H{0.34},C{0.16},H{0.50}}
  \pchead{רכיב & ערך & אופן חישוב ומעמד}
  \pcrow{משקלי דגם \latin{8B}, \latin{4-bit} & \num{4.0} \latin{GB} & \mth{8\cdot10^9\times4/8}; חישוב}
  \pcrow{זיכרון הקשר (\latin{8,192} טוקנים) & \num{1.1} \latin{GB} & \mth{2\cdot32\cdot8\cdot128\cdot2\cdot8192}; חישוב לפי ארכיטקטורה אופיינית}
  \pcrow{מודל ייצוג + מדרג & \num{1.3} \latin{GB} & הערכה; ייבדק בהתקנה}
  \pcrow{תקורת הפעלה & \num{1.0} \latin{GB} & הערכה}
  \pcrow{תצוגת \latin{3D} & \num{2.0} \latin{GB} & הערכה; תלוי במנוע}
  \pcrow{\textbf{סה"כ} & \textbf{\num{9.4} \latin{GB}} & מחושב; גדל ל-\mth{\approx12.6} \latin{GB} בהקשר של \latin{32,768}}
  \pcrow{זמן תגובה לשאילתה & \num{9.3} \latin{s} & \mth{0.3+1.5+300/40}; בהנחת \latin{40} טוקנים לשנייה}
\end{pctable}
\end{table}

\section{כלים מהתואר שהופעלו בפרק}\label{sec:mot}
הפרק נשען על שיטות מתואר ניהול הטכנולוגיות, והטבלה מפרטת איפה כל אחת הופעלה (טבלה~\ref{tab:mot}).
\begin{table}[!htbp]
\centering
\caption{שיטות מתואר ניהול הטכנולוגיות והיישום שלהן בפרק}\label{tab:mot}
\begin{pctable}{H{0.27},H{0.36},H{0.37}}
  \pchead{תחום בתואר & שיטה & יישום בפרק}
  \pcrow{ניתוח החלטות רב-קריטריוני \latin{(MCDA)} & מטריצת \latin{Pugh}, משקולות וניתוח רגישות & טבלאות~\ref{tab:pugh}--\ref{tab:sens}; נקודות מעבר בין חלופות}
  \pcrow{הנדסת מערכות \latin{(Systems Engineering)} & דרישות עם עקיבות, \latin{DFD}, ממשקים, מכונות מצבים & איורים~\ref{fig:dfd}--\ref{fig:sim}, טבלה~\ref{tab:schema}}
  \pcrow{חקר ביצועים \latin{(Operations Research)} & תורי \latin{M/M/1}, תכנון קיבולת & קיבולת המומחה; הגבלת קצב ההפקה}
  \pcrow{ניהול סיכונים & \mth{R=P\times I}, שערי \latin{Go/No-Go} & שער הבחירה; פרק~\ref{ch:risk}}
  \pcrow{ניהול פרויקטים & \latin{WBS}, תלויות ונתיב קריטי & הפיכת הארכיטקטורה לחבילות עבודה (פרק~\ref{ch:plan})}
\end{pctable}
\end{table}

\section{חדשנות}\label{sec:innovation}
\begin{pcbox}[pcGood]{מה חדש}
\textbf{ההצהרה:} שרשרת ניהולית-הנדסית מדידה ממסמך ולקח, דרך אחזור ויצירת תוצר מאושר, ועד תבנית \threeD{} שמאוכלסת אוטומטית מתצורת
\latin{JSON} ומתעדכנת עם שינוי גרסה --- ולא יצירת טקסט אוטומטית בלבד. \textbf{סוג החדשנות:} שילוב ויישום.\par
\textbf{ביחס למאמר הדגל:} מוסיף עיגון לגרסה ולעמוד, אישור מומחה לפני פרסום, מדידת זמן מול בסיס ידני ותרגול \latin{3D}
\cite{hua2025civask}. \textbf{ביחס למוצרי השוק:} פריסה מקומית מבודדת עם עקיבות מקור-גרסה-יחידה, שאינה מתועדת במוצרים
שנסקרו (טבלה~\ref{tab:products}).\par
\textbf{מה אינו חדש:} טכניקת ה-\latin{RAG} עצמה ומדדי הנאמנות נלקחים מן הספרות \cite{lewis2020retrieval,es2024ragas}.
\textbf{מענה לפער:} השילוב שהוגדר בסעיף~\ref{sec:gap}.
\textbf{הוכחה:} ניסוי משימות שקולות (ידני מול אב-טיפוס) ותרגיל \latin{2D} מול \latin{3D} בעל תוכן זהה, לפי טבלה~\ref{tab:eval}.
\end{pcbox}
\pcChapterEnd{alt}

\pcIfSolution{afteralternatives}
